---------------------------------
TODO tutorial
---------------------------------

\qquad \brandName realizó las siguientes acciones para comprometer completamente el dominio \enterpriseName.LOCAL.
1. El evaluador utilizó la herramienta Responder para obtener un hash de contraseña NTLMv2 de un usuario del dominio, bsmith.
2. Este hash de contraseña fue descifrado exitosamente fuera de línea usando la herramienta Hashcat, revelando la contraseña en texto claro del usuario, lo que otorgó un punto de acceso al dominio \enterpriseName.LOCAL, pero sin más privilegios que los de un usuario estándar del dominio.
3. Luego, el evaluador ejecutó BloodHound.py, una versión en Python de la popular herramienta de recolección SharpHound, para enumerar el dominio y crear representaciones visuales de las rutas de ataque. Tras una revisión, el evaluador encontró que existían múltiples usuarios privilegiados en el dominio configurados con Service Principal Names (SPNs), los cuales pueden ser aprovechados para realizar un ataque de Kerberoasting y recuperar tickets TGS de Kerberos para las cuentas. Estos tickets pueden ser descifrados fuera de línea usando Hashcat si se ha configurado una contraseña débil. Desde este punto, el evaluador utilizó la herramienta GetUserSPNs.py para realizar un ataque Kerberoasting dirigido contra la cuenta mssqlsvc, encontrando que dicha cuenta tenía derechos de administrador local sobre el host SQL01.\enterpriseName.LOCAL, lo que resultaba un objetivo interesante en el dominio.
4. El evaluador logró descifrar exitosamente la contraseña de esta cuenta fuera de línea, revelando su valor en texto claro.
5. El evaluador pudo autenticarse en el host SQL01.\enterpriseName.LOCAL y recuperar una contraseña en texto claro del registro del host al descifrar secretos LSA para una cuenta (srvadmin) que estaba configurada para inicio de sesión automático.
6. Esta cuenta srvadmin tenía derechos de administrador local sobre todos los servidores (excepto los Controladores de Dominio) en el dominio, por lo que el evaluador pudo iniciar sesión en el host MS01.\enterpriseName.LOCAL y recuperar un ticket TGT de Kerberos para un usuario conectado, pramirez, quien formaba parte del grupo de Administradores de Servidores de Nivel I, otorgándole al evaluador derechos de DCSync sobre el objeto del dominio. Este ataque puede ser utilizado para recuperar el hash de contraseña NTLM de cualquier usuario en el dominio, resultando en el compromiso del dominio y persistencia mediante un Golden Ticket.
7. El evaluador utilizó la herramienta Rubeus para extraer el ticket TGT de Kerberos del usuario pramirez y realizar un ataque de Pass-the-Ticket para autenticarse como este usuario.
8. Finalmente, el evaluador logró realizar un ataque DCSync después de autenticarse exitosamente con esta cuenta de usuario utilizando la herramienta Mimikatz, lo que resultó en el compromiso total del dominio.
